\section{Lorentz Invariance on Every Stratum}
\begin{quote}\small
\noindent\textbf{Status.} \emph{Legacy physical-readout block inside the active main body.}\par
\noindent\textbf{Block function.} This section isolates Lorentz invariance as a stratumwise consequence of the structural core before later physical applications are read off from it.
\end{quote}
\label{sec:lorentz}
%% ============================================================

\begin{theorem}[Stratified Lorentz invariance]\label{thm:lorentz}
On every stratum $S_n$, the quadratic invariant
\[
|q_n|^2=\tau_n^2
\]
is preserved under every HC-admissible frame transformation induced by
the global Heisenberg-compensator matrix.
In particular, for the involutive frame transformation
\[
J:\ z^- \mapsto -z^-,
\]
the proper-time invariant $\tau_n$ is unchanged.
Moreover, every admissible relabeling of the frame data that preserves
the same involution $J$ leaves both $\tau_n$ and the compensator
\[
C(f)=\tfrac12(f+Jf)
\]
unchanged.
\end{theorem}

\begin{proof}
By the general MM-tools, a change of admissible realization is not a
change of structural object.
By the HC-tools, a frame transformation is physically admissible only if
it preserves the same compensator structure, i.e.\ the same involution
$J$.

The transformation $J:z^-\mapsto -z^-$ changes only the frame-sensitive
sign, not the underlying compensated structural content.
Hence the quadratic invariant associated with the stratum remains the
same, so
\[
|q_n|^2=\tau_n^2
\]
is preserved.

Likewise, the compensator
\[
C=\tfrac12(\mathrm{id}+J)
\]
depends only on the involution $J$ itself.
Therefore any HC-admissible relabeling that preserves the same
$J$-structure leaves both $\tau_n$ and $C$ unchanged.
\end{proof}

%% ============================================================

